% Adaptive Time-Stepping
% Source: docs/adaptive_timestepping.md (migrated to LaTeX)
% Uses macros from preamble.tex

\section{Adaptive Time-Stepping}
\label{sec:adaptive-timestepping}

\textbf{Status:} Implemented (CVODE + scipy + IDA + leapfrog) \\
\textbf{Date:} 2026-02-23

Tolerance-controlled adaptive time-stepping replaces manual $\Delta t$ selection
with error-controlled solvers that automatically choose step sizes. Users
specify accuracy targets (via \texttt{--rtol}/\texttt{--atol}) instead of guessing a
stable $\Delta t$.

\subsection{Problem Statement}
\label{sec:at-problem}

Fixed-step integrators (leapfrog) require the user to manually satisfy the
CFL condition: $\Delta t < \Delta x / c_\mathrm{max}$. This has three problems:

\begin{enumerate}
  \item \textbf{Fragility.} A $\Delta t$ that works for one parameter set can be unstable
    for another. Getting it wrong produces garbage output that can look
    plausible.

  \item \textbf{Inefficiency.} A $\Delta t$ chosen for the fastest mode wastes compute on
    slow phases. Resonant mixing (the Gertsenshtein effect) has beat
    frequencies much slower than the carrier --- fixed steps resolve every
    carrier cycle even when nothing interesting is happening. Berlin et al.\
    (2024)~\cite{berlin2024axionphoton} document these numerical challenges.

  \item \textbf{No accuracy guarantee.} Leapfrog energy error is $O(\Delta t^2)$ --- the
    shadow Hamiltonian drifts by an amount proportional to $\Delta t^2$. There is no
    built-in error control to bound this.
\end{enumerate}

\subsection{Four Solver Paths}
\label{sec:solver-paths}

\subsubsection{\texttt{--scheme auto} (Default): Intelligent Selection}
\label{sec:scheme-auto}

Auto-selection \textbf{always} picks a tolerance-controlled adaptive solver.
Leapfrog is opt-in only via \texttt{--scheme leapfrog}.

Detection algorithm:

\begin{table}[htbp]
\centering
\caption{Auto-selection priority for solver schemes.}
\label{tab:auto-selection}
\begin{tabular}{@{}llll@{}}
\toprule
Priority & Condition & Solver & Reason \\
\midrule
1 & Modal eligible (flat metric + periodic BCs + time-independent + supported operators) & \textbf{Modal} & Machine-precision eigendecomposition, $O(1)$ in $t_\mathrm{end}$. See \cref{sec:scheme-modal} \\
2 & Constraint equations (\texttt{time\_order=0}) not modal-eligible & IDA & Algebraic constraints need DAE residual form \\
3 & First-order equations (\texttt{time\_order=1}) & IDA & Diffusion/transport need implicit handling \\
4 & Dissipation (\texttt{first\_derivative\_t} in RHS) & IDA & Friction breaks symplecticity, BDF handles well \\
5 & No canonical Hamiltonian & CVODE (warning) & Missing canonical section --- likely hand-crafted JSON \\
6 & Pure wave (2nd-order, Hamiltonian) & CVODE & Adaptive BDF, no DAE overhead \\
\bottomrule
\end{tabular}
\end{table}

\subsubsection{\texttt{--scheme cvode}: SUNDIALS CVODE (Adaptive ODE)}
\label{sec:scheme-cvode}

Pure ODE solver using SUNDIALS CVODE via scikit-sundae. Supports BDF
(stiff, order 1--5) and Adams (non-stiff, order 1--12).

\begin{lstlisting}[style=bash]
tidal simulate spec.json --scheme cvode                 # BDF default
tidal simulate spec.json --scheme cvode --method Adams   # Non-stiff
tidal simulate spec.json --scheme cvode --rtol 1e-6      # Custom tolerance
\end{lstlisting}

\textbf{Best for:} Wave systems (the auto-selected default for pure wave
equations). BDF handles the mild stiffness from spatial discretization
(eigenvalues scale with $1/\Delta x^2$). In-place RHS evaluation avoids
per-call allocation.

Reference: Hindmarsh et al.~\cite{hindmarsh2005sundials}, ``SUNDIALS: Suite of Nonlinear and
Differential/Algebraic Equation Solvers'', ACM TOMS 31(3).

\subsubsection{\texttt{--scheme modal}: Fourier Modal Solver (Exact Eigendecomposition)}
\label{sec:scheme-modal}

Transforms spatial grid to Fourier space, builds per-mode evolution matrix,
and eigendecomposes for exact $y(t) = \exp(A \cdot t) \cdot y_0$ solutions. No time-stepping,
no CFL condition, no tolerance parameter needed.

\begin{lstlisting}[style=bash]
tidal simulate spec.json --scheme modal                  # Auto for eligible
\end{lstlisting}

\textbf{Best for:} Periodic-domain wave systems with time-independent coefficients
(the auto-selected default for eligible systems). Machine-precision accuracy
($\sim 10^{-16}$ error). Cost is $O(1)$ in simulation time --- eigendecomposition is done
once, evaluation at any $t$ is just $\exp(\lambda \cdot \Delta t)$. Speedup over CVODE: 57--1451$\times$
on \texttt{coupled\_scalars} ($N=64$--$256$), 3.4--7.7$\times$ on Gertsenshtein ($N=256$, $t=50$--$500$).
Constraint systems (e.g., Proca $A_0$) are handled via Fourier Schur complement
elimination --- 9.5$\times$ speedup over IDA on \texttt{coupled\_proca\_3d}.

\textbf{Requires:} Flat metric, all-periodic BCs, time-independent coefficients,
supported Fourier operators. Position-dependent coefficients are handled via
$k$-space convolution.

\subsubsection{\texttt{--scheme ida}: SUNDIALS IDA (Adaptive DAE)}
\label{sec:scheme-ida}

DAE solver using SUNDIALS IDA. Handles algebraic constraints via residual
form $F(t, y, y') = 0$.

\begin{lstlisting}[style=bash]
tidal simulate spec.json --scheme ida                    # Auto for DAE
tidal simulate spec.json --scheme ida --rtol 1e-6        # Custom tolerance
\end{lstlisting}

\textbf{Best for:} Constraint systems (Chern--Simons, coupled Proca), dissipative
systems (Hubble friction), and any system with mixed time-derivative orders.

\subsubsection{\texttt{--scheme scipy}: scipy \texttt{solve\_ivp} (Adaptive ODE)}
\label{sec:scheme-scipy}

Wraps \texttt{scipy.integrate.solve\_ivp}. Default method is DOP853 (8th-order
Dormand--Prince), an explicit one-step method with excellent accuracy per
function evaluation.

\begin{lstlisting}[style=bash]
tidal simulate spec.json --scheme scipy                  # DOP853 default
tidal simulate spec.json --scheme scipy --method RK45    # Lower-order explicit
tidal simulate spec.json --scheme scipy --method Radau   # Implicit (stiff)
tidal simulate spec.json --scheme scipy --method BDF     # Implicit multi-step
\end{lstlisting}

\textbf{Best for:} Smooth non-stiff wave problems where explicit high-order
methods outperform implicit BDF (fewer function evaluations per accuracy
unit). Also useful when scipy's Radau is preferred over SUNDIALS.

Reference: Dormand \& Prince (1980), ``A family of embedded Runge--Kutta
formulae'', J.\ Comp.\ Appl.\ Math.\ 6, 19--26.

\subsubsection{\texttt{--scheme leapfrog}: Symplectic (Opt-In, Fixed $\Delta t$)}
\label{sec:scheme-leapfrog}

St\"ormer--Verlet symplectic integrator. Preserves a shadow Hamiltonian to
machine precision --- zero secular energy drift.

\begin{lstlisting}[style=bash]
tidal simulate spec.json --scheme leapfrog               # CFL-auto dt
tidal simulate spec.json --scheme leapfrog --dt 0.01     # Manual dt
\end{lstlisting}

\textbf{Best for:} Long-time Hamiltonian evolution where zero energy drift
matters more than per-step accuracy. Not suitable for dissipative systems,
constraints, or when tolerance control is needed.

Reference: Hairer, Lubich \& Wanner (2006), ``Geometric Numerical
Integration'', Springer, Chapter~VI.

\subsubsection{Choosing a Solver}
\label{sec:choosing-solver}

\begin{table}[htbp]
\centering
\caption{Solver recommendation by scenario.}
\label{tab:solver-recommendation}
\begin{tabular}{@{}ll@{}}
\toprule
Scenario & Recommended \\
\midrule
Default (any system) & \texttt{--scheme auto} (always adaptive) \\
Wave equation, production accuracy & \texttt{--scheme cvode --rtol 1e-8} \\
Constraint system ($A_0$, gauge fields) & \texttt{--scheme ida} (auto-selected) \\
Dissipative system (friction, damping) & \texttt{--scheme ida} (auto-selected) \\
Smooth non-stiff, high accuracy & \texttt{--scheme scipy --rtol 1e-10} \\
Long-time energy conservation & \texttt{--scheme leapfrog} \\
Stiff system ($m^2 \Delta x^2 \gg 1$) & \texttt{--scheme scipy --method Radau} \\
Convergence study & \texttt{--scheme cvode --rtol 1e-12 --atol 1e-14} \\
\bottomrule
\end{tabular}
\end{table}

\subsection{CLI Reference}
\label{sec:cli-reference}

\begin{table}[htbp]
\centering
\caption{CLI flags for adaptive time-stepping.}
\label{tab:cli-flags}
\begin{tabular}{@{}llll@{}}
\toprule
Flag & Type & Default & Description \\
\midrule
\texttt{--scheme} & \texttt{\{auto,ida,leapfrog,cvode,scipy\}} & \texttt{auto} & Solver selection \\
\texttt{--rtol} & float & \texttt{1e-8} & Relative tolerance (adaptive solvers) \\
\texttt{--atol} & float & \texttt{1e-10} & Absolute tolerance (adaptive solvers) \\
\texttt{--method} & str & auto & Integration method (cvode: BDF/Adams; scipy: DOP853/RK45/Radau/BDF) \\
\texttt{--max-step} & float & auto & Maximum step size \\
\texttt{--dt} & float & CFL & Fixed time step (leapfrog only) \\
\bottomrule
\end{tabular}
\end{table}

\subsubsection{Default Tolerances}
\label{sec:default-tolerances}

All adaptive solvers use \texttt{rtol=1e-8}, \texttt{atol=1e-10} by default. These are
production-quality defaults that bound the temporal integration error well
below spatial discretization error for typical grid resolutions (32--256
points).

\subsubsection{Validation Rules}
\label{sec:validation-rules}

\begin{itemize}
  \item \texttt{--rtol} must be positive
  \item \texttt{--atol} must be positive
  \item \texttt{--max-step} must be positive
  \item \texttt{--method} is only meaningful with \texttt{--scheme cvode} or \texttt{--scheme scipy}
\end{itemize}

\subsection{Architecture}
\label{sec:at-architecture}

\subsubsection{Shared Infrastructure}
\label{sec:shared-infrastructure}

All four solvers share core infrastructure, minimizing code duplication:

\begin{itemize}
  \item \textbf{\texttt{compute\_force()}} (\texttt{leapfrog.py}): Evaluates $\mathrm{d}\pi/\mathrm{d}t = F(q)$ for all
    momentum slots. Used by leapfrog, CVODE, and scipy.
  \item \textbf{\texttt{compute\_velocity()}} (\texttt{leapfrog.py}): Evaluates $\mathrm{d}q/\mathrm{d}t = K^{-1}(\pi - S)$
    for all field slots. Used by leapfrog, CVODE, and scipy.
  \item \textbf{\texttt{RHSEvaluator}} (\texttt{rhs.py}): Unified operator + coefficient application.
    Used by all four solvers.
  \item \textbf{\texttt{CoefficientEvaluator}} (\texttt{coefficients.py}): 4-level cache (L0--L3).
    Used by all adaptive solvers.
  \item \textbf{\texttt{build\_jacobian\_sparsity()}} (\texttt{sparsity.py}): Jacobian sparsity pattern
    for implicit solvers. Used by IDA, CVODE (sparse linsolver), and scipy
    (Radau/BDF).
  \item \textbf{\texttt{StateLayout}}, \textbf{\texttt{FieldSet}}, \textbf{\texttt{GridInfo}}: State management shared
    by all solvers.
\end{itemize}

\subsubsection{Result Format}
\label{sec:result-format}

All solvers return the same dictionary:

\begin{lstlisting}[style=python]
{
    "t": np.ndarray,     # (num_snapshots,) time points
    "y": np.ndarray,     # (num_snapshots, total_size) state snapshots
    "success": bool,     # True if integration completed
    "message": str,      # Human-readable status
}
\end{lstlisting}

\subsubsection{Linear Solver Selection (SUNDIALS)}
\label{sec:linear-solver}

Both CVODE and IDA use the same thresholds for automatic linear solver
selection based on system size:

\begin{table}[htbp]
\centering
\caption{Automatic linear solver selection by system size.}
\label{tab:linear-solver}
\begin{tabular}{@{}lll@{}}
\toprule
System Size & Linear Solver & Rationale \\
\midrule
$N \leq 2{,}000$ & Dense & Direct solve, minimal overhead \\
$2{,}000 < N \leq 200{,}000$ & Sparse (SuperLU) & Exploits sparsity pattern \\
$N > 200{,}000$ & GMRES (iterative) & Avoids $O(N^2)$ memory \\
\bottomrule
\end{tabular}
\end{table}

\subsubsection{Sparsity Reuse}
\label{sec:sparsity-reuse}

The IDA sparsity builder (\texttt{build\_jacobian\_sparsity()}) produces a pattern
for $J = \mathrm{d}F/\mathrm{d}y + c_j \cdot \mathrm{d}F/\mathrm{d}y'$. For CVODE (ODE form), the needed pattern
is $J = \mathrm{d}f/\mathrm{d}y$. The IDA pattern is a \textbf{superset} of CVODE's true pattern:
extra diagonal entries from $c_j \cdot I$ are always nonzero anyway. This means
the same sparsity builder works for both solvers --- a key simplification.

For scipy's implicit methods (Radau, BDF), the same pattern is passed as
\texttt{jac\_sparsity} to \texttt{solve\_ivp}.

\subsection{Design Decisions}
\label{sec:design-decisions}

\subsubsection{Why CVODE as Default for Wave Systems (Not Leapfrog)?}
\label{sec:why-cvode}

\begin{itemize}
  \item \textbf{Tolerance-controlled:} Energy error bounded by rtol, not $O(\Delta t^2)$
  \item \textbf{Adaptive:} Automatically picks step sizes, no CFL tuning
  \item \textbf{Safe default:} Works correctly for all wave systems out of the box
  \item \textbf{Stiffness handling:} BDF handles mild stiffness from $1/\Delta x^2$ eigenvalues
  \item Leapfrog remains available for users who need symplectic integration
\end{itemize}

\subsubsection{Why CVODE over scipy as Default?}
\label{sec:cvode-vs-scipy}

\begin{itemize}
  \item Method-of-lines PDE discretization is mildly stiff (eigenvalues near
    imaginary axis scale with $1/\Delta x^2$). BDF handles this robustly.
  \item DOP853 would need a CFL-like \texttt{max\_step} guard to prevent instability
    during quiescent phases.
  \item Same SUNDIALS ecosystem as IDA --- consistent tolerance semantics.
  \item In-place RHS (\texttt{rhsfn(t, y, yp)}) avoids allocation per call.
\end{itemize}

\subsubsection{Why Also Provide scipy?}
\label{sec:why-scipy}

\begin{itemize}
  \item DOP853 (8th-order explicit) excels for smooth non-stiff waves --- fewer
    function evaluations than BDF for the same accuracy.
  \item One-step method = no startup cost (relevant for short simulations).
  \item Good option when implicit BDF is overkill.
\end{itemize}

\subsubsection{Why Four Solvers?}
\label{sec:why-four}

Each has a unique, irreplaceable strength:

\begin{table}[htbp]
\centering
\caption{Unique capability of each solver backend.}
\label{tab:unique-capabilities}
\begin{tabular}{@{}ll@{}}
\toprule
Solver & Unique Capability \\
\midrule
Leapfrog & Symplectic --- zero secular energy drift \\
CVODE & Implicit adaptive ODE --- mildly stiff, tolerance-controlled \\
IDA & DAE --- algebraic constraints (irreplaceable for Chern--Simons) \\
scipy & Explicit adaptive --- DOP853 8th-order for smooth problems \\
\bottomrule
\end{tabular}
\end{table}

\subsection{Usage Examples}
\label{sec:usage-examples}

\begin{lstlisting}[style=bash]
# Default -- auto-selects best adaptive solver
tidal simulate spec.json

# Custom tolerances
tidal simulate spec.json --rtol 1e-6 --atol 1e-8

# CVODE with Adams (non-stiff, high-order)
tidal simulate spec.json --scheme cvode --method Adams

# High-accuracy explicit (DOP853)
tidal simulate spec.json --scheme scipy --rtol 1e-10

# IDA with custom tolerances
tidal simulate spec.json --scheme ida --rtol 1e-6

# Symplectic (opt-in, no tolerance control)
tidal simulate spec.json --scheme leapfrog
\end{lstlisting}

\subsection{References}
\label{sec:at-references}

\begin{itemize}
  \item \textbf{Hindmarsh et al.\ (2005)}~\cite{hindmarsh2005sundials}, ``SUNDIALS: Suite of Nonlinear and
    Differential/Algebraic Equation Solvers'', ACM TOMS 31(3). --- CVODE and
    IDA algorithms, BDF/Adams methods, SUNDIALS architecture.

  \item \textbf{Hairer, N\o rsett \& Wanner (1993)}, \emph{Solving Ordinary Differential
    Equations I: Nonstiff Problems}, Springer. --- DOP853, embedded Runge--Kutta
    theory, error estimation.

  \item \textbf{Hairer \& Wanner (1996)}, \emph{Solving Ordinary Differential Equations II:
    Stiff and Differential-Algebraic Problems}, Springer. --- BDF, Radau,
    stiffness theory, DAE index theory.

  \item \textbf{Dormand \& Prince (1980)}, ``A family of embedded Runge--Kutta formulae'',
    J.\ Comp.\ Appl.\ Math.\ 6, 19--26. --- The RK method family underlying DOP853.

  \item \textbf{Hairer, Lubich \& Wanner (2006)}, \emph{Geometric Numerical Integration},
    Springer. --- Symplectic integrators (St\"ormer--Verlet/leapfrog).

  \item \textbf{Berlin et al.\ (2024)}~\cite{berlin2024axionphoton}, ``Numerical analysis of resonant axion-photon
    mixing'', arXiv:2405.08865. --- Numerical challenges in resonant wave mixing
    analogous to TIDAL's Gertsenshtein simulation.

  \item \textbf{Courant, Friedrichs \& Lewy (1928)}, ``\"Uber die partiellen
    Differenzengleichungen der mathematischen Physik'', Math.\ Ann.\ 100, 32--74.
    --- The CFL stability condition.

  \item \textbf{scikit-sundae} (NREL, BSD-3 license) --- Python wrapper for SUNDIALS
    CVODE and IDA solvers.
\end{itemize}
